% Options for packages loaded elsewhere
\PassOptionsToPackage{unicode}{hyperref}
\PassOptionsToPackage{hyphens}{url}
\documentclass[
]{article}
\usepackage{xcolor}
\usepackage[margin=1in]{geometry}
\usepackage{amsmath,amssymb}
\setcounter{secnumdepth}{-\maxdimen} % remove section numbering
\usepackage{iftex}
\ifPDFTeX
  \usepackage[T1]{fontenc}
  \usepackage[utf8]{inputenc}
  \usepackage{textcomp} % provide euro and other symbols
\else % if luatex or xetex
  \usepackage{unicode-math} % this also loads fontspec
  \defaultfontfeatures{Scale=MatchLowercase}
  \defaultfontfeatures[\rmfamily]{Ligatures=TeX,Scale=1}
\fi
\usepackage{lmodern}
\ifPDFTeX\else
  % xetex/luatex font selection
\fi
% Use upquote if available, for straight quotes in verbatim environments
\IfFileExists{upquote.sty}{\usepackage{upquote}}{}
\IfFileExists{microtype.sty}{% use microtype if available
  \usepackage[]{microtype}
  \UseMicrotypeSet[protrusion]{basicmath} % disable protrusion for tt fonts
}{}
\makeatletter
\@ifundefined{KOMAClassName}{% if non-KOMA class
  \IfFileExists{parskip.sty}{%
    \usepackage{parskip}
  }{% else
    \setlength{\parindent}{0pt}
    \setlength{\parskip}{6pt plus 2pt minus 1pt}}
}{% if KOMA class
  \KOMAoptions{parskip=half}}
\makeatother
\usepackage{color}
\usepackage{fancyvrb}
\newcommand{\VerbBar}{|}
\newcommand{\VERB}{\Verb[commandchars=\\\{\}]}
\DefineVerbatimEnvironment{Highlighting}{Verbatim}{commandchars=\\\{\}}
% Add ',fontsize=\small' for more characters per line
\usepackage{framed}
\definecolor{shadecolor}{RGB}{248,248,248}
\newenvironment{Shaded}{\begin{snugshade}}{\end{snugshade}}
\newcommand{\AlertTok}[1]{\textcolor[rgb]{0.94,0.16,0.16}{#1}}
\newcommand{\AnnotationTok}[1]{\textcolor[rgb]{0.56,0.35,0.01}{\textbf{\textit{#1}}}}
\newcommand{\AttributeTok}[1]{\textcolor[rgb]{0.13,0.29,0.53}{#1}}
\newcommand{\BaseNTok}[1]{\textcolor[rgb]{0.00,0.00,0.81}{#1}}
\newcommand{\BuiltInTok}[1]{#1}
\newcommand{\CharTok}[1]{\textcolor[rgb]{0.31,0.60,0.02}{#1}}
\newcommand{\CommentTok}[1]{\textcolor[rgb]{0.56,0.35,0.01}{\textit{#1}}}
\newcommand{\CommentVarTok}[1]{\textcolor[rgb]{0.56,0.35,0.01}{\textbf{\textit{#1}}}}
\newcommand{\ConstantTok}[1]{\textcolor[rgb]{0.56,0.35,0.01}{#1}}
\newcommand{\ControlFlowTok}[1]{\textcolor[rgb]{0.13,0.29,0.53}{\textbf{#1}}}
\newcommand{\DataTypeTok}[1]{\textcolor[rgb]{0.13,0.29,0.53}{#1}}
\newcommand{\DecValTok}[1]{\textcolor[rgb]{0.00,0.00,0.81}{#1}}
\newcommand{\DocumentationTok}[1]{\textcolor[rgb]{0.56,0.35,0.01}{\textbf{\textit{#1}}}}
\newcommand{\ErrorTok}[1]{\textcolor[rgb]{0.64,0.00,0.00}{\textbf{#1}}}
\newcommand{\ExtensionTok}[1]{#1}
\newcommand{\FloatTok}[1]{\textcolor[rgb]{0.00,0.00,0.81}{#1}}
\newcommand{\FunctionTok}[1]{\textcolor[rgb]{0.13,0.29,0.53}{\textbf{#1}}}
\newcommand{\ImportTok}[1]{#1}
\newcommand{\InformationTok}[1]{\textcolor[rgb]{0.56,0.35,0.01}{\textbf{\textit{#1}}}}
\newcommand{\KeywordTok}[1]{\textcolor[rgb]{0.13,0.29,0.53}{\textbf{#1}}}
\newcommand{\NormalTok}[1]{#1}
\newcommand{\OperatorTok}[1]{\textcolor[rgb]{0.81,0.36,0.00}{\textbf{#1}}}
\newcommand{\OtherTok}[1]{\textcolor[rgb]{0.56,0.35,0.01}{#1}}
\newcommand{\PreprocessorTok}[1]{\textcolor[rgb]{0.56,0.35,0.01}{\textit{#1}}}
\newcommand{\RegionMarkerTok}[1]{#1}
\newcommand{\SpecialCharTok}[1]{\textcolor[rgb]{0.81,0.36,0.00}{\textbf{#1}}}
\newcommand{\SpecialStringTok}[1]{\textcolor[rgb]{0.31,0.60,0.02}{#1}}
\newcommand{\StringTok}[1]{\textcolor[rgb]{0.31,0.60,0.02}{#1}}
\newcommand{\VariableTok}[1]{\textcolor[rgb]{0.00,0.00,0.00}{#1}}
\newcommand{\VerbatimStringTok}[1]{\textcolor[rgb]{0.31,0.60,0.02}{#1}}
\newcommand{\WarningTok}[1]{\textcolor[rgb]{0.56,0.35,0.01}{\textbf{\textit{#1}}}}
\usepackage{graphicx}
\makeatletter
\newsavebox\pandoc@box
\newcommand*\pandocbounded[1]{% scales image to fit in text height/width
  \sbox\pandoc@box{#1}%
  \Gscale@div\@tempa{\textheight}{\dimexpr\ht\pandoc@box+\dp\pandoc@box\relax}%
  \Gscale@div\@tempb{\linewidth}{\wd\pandoc@box}%
  \ifdim\@tempb\p@<\@tempa\p@\let\@tempa\@tempb\fi% select the smaller of both
  \ifdim\@tempa\p@<\p@\scalebox{\@tempa}{\usebox\pandoc@box}%
  \else\usebox{\pandoc@box}%
  \fi%
}
% Set default figure placement to htbp
\def\fps@figure{htbp}
\makeatother
\setlength{\emergencystretch}{3em} % prevent overfull lines
\providecommand{\tightlist}{%
  \setlength{\itemsep}{0pt}\setlength{\parskip}{0pt}}
\usepackage{bookmark}
\IfFileExists{xurl.sty}{\usepackage{xurl}}{} % add URL line breaks if available
\urlstyle{same}
% fallback for those not using the hyperref driver hyperxmp:
\makeatletter
\@ifundefined{xmpquote}{\newcommand{\xmpquote}[1]{#1}}{}
\makeatother
\hypersetup{
  pdftitle={Actividad 1.1 Normalidad univariada},
  pdfauthor={Lucero Díaz},
  hidelinks,
  pdfcreator={LaTeX via pandoc}}

\title{Actividad 1.1 Normalidad univariada}
\author{Lucero Díaz}
\date{2026-09-22}

\begin{document}
\maketitle

\subsection{Sesgo y curtosis en la Normalidad y fuera de
ella}\label{sesgo-y-curtosis-en-la-normalidad-y-fuera-de-ella}

\subsubsection{1. En la normalidad}\label{en-la-normalidad}

\begin{Shaded}
\begin{Highlighting}[]
\CommentTok{\# A. Dibuja el gráfico de la distribución normal con media cero y desviación estándar 100}

\NormalTok{Pob\_normal }\OtherTok{=} \FunctionTok{dnorm}\NormalTok{(}\SpecialCharTok{{-}}\DecValTok{400}\SpecialCharTok{:}\DecValTok{400}\NormalTok{, }\DecValTok{0}\NormalTok{, }\DecValTok{100}\NormalTok{)}
\FunctionTok{plot}\NormalTok{(}\SpecialCharTok{{-}}\DecValTok{400}\SpecialCharTok{:}\DecValTok{400}\NormalTok{, Pob\_normal, }\AttributeTok{type=}\StringTok{"l"}\NormalTok{, }\AttributeTok{main =} \StringTok{"Distribución Normal media=0, sigma=100"}\NormalTok{)}
\end{Highlighting}
\end{Shaded}

\pandocbounded{\includegraphics[keepaspectratio]{Actividad-1_files/figure-latex/unnamed-chunk-1-1.pdf}}

\begin{Shaded}
\begin{Highlighting}[]
\CommentTok{\# B. Genera 150 datos aleatorios provenientes de una distribución normal con media cero y desviación estándar 100 y dibuja el histograma de los datos generados.}
\NormalTok{x1}\OtherTok{=}\FunctionTok{rnorm}\NormalTok{(}\DecValTok{100}\NormalTok{,}\DecValTok{0}\NormalTok{,}\DecValTok{100}\NormalTok{)}
\FunctionTok{hist}\NormalTok{(x1, }\AttributeTok{col=}\DecValTok{0}\NormalTok{, }\AttributeTok{main =} \StringTok{"Distribución empírica de Normal media=0, sigma=100"}\NormalTok{, }\AttributeTok{freq=}\ConstantTok{FALSE}\NormalTok{)}
\end{Highlighting}
\end{Shaded}

\pandocbounded{\includegraphics[keepaspectratio]{Actividad-1_files/figure-latex/unnamed-chunk-2-1.pdf}}

\begin{Shaded}
\begin{Highlighting}[]
\CommentTok{\# C. Calcula el sesgo y la curtosis de los datos generados:}
\FunctionTok{library}\NormalTok{ (moments)}
\NormalTok{sg1 }\OtherTok{=} \FunctionTok{skewness}\NormalTok{(x1)}
\NormalTok{k1 }\OtherTok{=} \FunctionTok{kurtosis}\NormalTok{(x1)}

\FunctionTok{cat}\NormalTok{(}\StringTok{"Sesgo:"}\NormalTok{, sg1, }\StringTok{"}\SpecialCharTok{\textbackslash{}n}\StringTok{"}\NormalTok{)}
\end{Highlighting}
\end{Shaded}

\begin{verbatim}
## Sesgo: 0.3385956
\end{verbatim}

\begin{Shaded}
\begin{Highlighting}[]
\FunctionTok{cat}\NormalTok{(}\StringTok{"Curtosis:"}\NormalTok{, k1)}
\end{Highlighting}
\end{Shaded}

\begin{verbatim}
## Curtosis: 2.730511
\end{verbatim}

\begin{Shaded}
\begin{Highlighting}[]
\CommentTok{\# D. Analiza la normalidad de los datos generados:}
\FunctionTok{jarque.test}\NormalTok{(x1)}
\end{Highlighting}
\end{Shaded}

\begin{verbatim}
## 
##  Jarque-Bera Normality Test
## 
## data:  x1
## JB = 2.2134, p-value = 0.3307
## alternative hypothesis: greater
\end{verbatim}

\begin{Shaded}
\begin{Highlighting}[]
\FunctionTok{qqnorm}\NormalTok{(x1, }\AttributeTok{main =} \StringTok{"QQ{-}Plot de datos Normales"}\NormalTok{)}
\FunctionTok{qqline}\NormalTok{(x1, }\AttributeTok{col =} \StringTok{"red"}\NormalTok{)}
\end{Highlighting}
\end{Shaded}

\pandocbounded{\includegraphics[keepaspectratio]{Actividad-1_files/figure-latex/unnamed-chunk-4-1.pdf}}

\begin{itemize}
\tightlist
\item
  Compara el gráfico de la distribución teórica (punto A) con el gráfico
  de la distribución empírica (punto B), ¿en qué son iguales y en qué
  son diferentes?
\end{itemize}

Ambas gráficas se muestran centradas en 0, así como también un rango de
variabilidad similar. En donde mayormente difieren sería en la
variabilidad de cada barra, lo que termina afectando en la distribución
empírica.

\begin{itemize}
\tightlist
\item
  Relaciona la distribución empírica (punto B) con los valores de sesgo
  y curtosis (punto C) y el QQPlot y la prueba de normalidad (punto D)
\end{itemize}

Para analizar la distribución empírica mostrada, veremos cada punto paso
a paso (C y D). En el caso del punto C, se nos muestran tanto el sesgo
como la curtosis muy cercanos a 0, lo que indica simetría y una
distribución mesocúrtica; mientras que en el punto D, la QQ-Plot nos
muestra como la mayoría de los datos si siguen la línea teórica, a pesar
de que ambas colas tengan algunos datos fuera de la línea. Finalmente,
la prueba de normalidad Jarque-Bera nos arroja un p-value mayor a 0.05,
lo que confirma la hipótesis nula y con esto tenemos evidencia de que la
distribución empírica si puede considerarse una distribución normal.

\subsubsection{2. Fuera de la normalidad}\label{fuera-de-la-normalidad}

\begin{Shaded}
\begin{Highlighting}[]
\CommentTok{\# Configuración del panel para mostrar 3 gráficos}
\FunctionTok{par}\NormalTok{(}\AttributeTok{mfrow =} \FunctionTok{c}\NormalTok{(}\DecValTok{1}\NormalTok{, }\DecValTok{3}\NormalTok{))}

\CommentTok{\# Crear el vector del dominio x (de 0 a 10 con incrementos finos para una curva suave)}
\NormalTok{x11 }\OtherTok{=} \FunctionTok{seq}\NormalTok{(}\DecValTok{0}\NormalTok{, }\DecValTok{10}\NormalTok{, }\AttributeTok{by =} \FloatTok{0.01}\NormalTok{)}

\CommentTok{\# Evaluar la densidad dgamma (o dexp)}
\NormalTok{Pob\_muy\_derecha }\OtherTok{=} \FunctionTok{dgamma}\NormalTok{(x11, }\AttributeTok{shape =} \DecValTok{1}\NormalTok{, }\AttributeTok{rate =} \DecValTok{1}\NormalTok{)}

\CommentTok{\# Graficar la línea}
\FunctionTok{plot}\NormalTok{(x11, Pob\_muy\_derecha, }\AttributeTok{type =} \StringTok{"l"}\NormalTok{, }
     \AttributeTok{main =} \StringTok{"Distribución Muy Sesgada a la Derecha (Gamma shape=1, rate=1)"}\NormalTok{,}
     \AttributeTok{xlab =} \StringTok{"x"}\NormalTok{, }\AttributeTok{ylab =} \StringTok{"Densidad"}\NormalTok{)}

\CommentTok{\# Definir el dominio x}
\NormalTok{x12 }\OtherTok{=} \FunctionTok{seq}\NormalTok{(}\DecValTok{0}\NormalTok{, }\DecValTok{5}\NormalTok{, }\AttributeTok{by =} \FloatTok{0.01}\NormalTok{)}

\CommentTok{\# Evaluar la densidad dweibull}
\NormalTok{Pob\_lig\_derecha }\OtherTok{=} \FunctionTok{dweibull}\NormalTok{(x12, }\AttributeTok{shape =} \DecValTok{2}\NormalTok{, }\AttributeTok{scale =} \DecValTok{1}\NormalTok{)}

\CommentTok{\# Graficar la línea}
\FunctionTok{plot}\NormalTok{(x12, Pob\_lig\_derecha, }\AttributeTok{type =} \StringTok{"l"}\NormalTok{, }
     \AttributeTok{main =} \StringTok{"Distribución Ligeramente Sesgada a la Derecha (Weibull shape=2, scale=1)"}\NormalTok{,}
     \AttributeTok{xlab =} \StringTok{"x"}\NormalTok{, }\AttributeTok{ylab =} \StringTok{"Densidad"}\NormalTok{)}

\CommentTok{\# Definir el dominio x}
\NormalTok{x13 }\OtherTok{=} \FunctionTok{seq}\NormalTok{(}\DecValTok{0}\NormalTok{, }\FloatTok{1.5}\NormalTok{, }\AttributeTok{by =} \FloatTok{0.01}\NormalTok{)}

\CommentTok{\# Evaluar la densidad dweibull}
\NormalTok{Pob\_izquierda }\OtherTok{=} \FunctionTok{dweibull}\NormalTok{(x13, }\AttributeTok{shape =} \DecValTok{10}\NormalTok{, }\AttributeTok{scale =} \DecValTok{1}\NormalTok{)}

\CommentTok{\# Graficar la línea}
\FunctionTok{plot}\NormalTok{(x13, Pob\_izquierda, }\AttributeTok{type =} \StringTok{"l"}\NormalTok{, }
     \AttributeTok{main =} \StringTok{"Distribución Sesgada a la Izquierda (Weibull shape=10, scale=1)"}\NormalTok{,}
     \AttributeTok{xlab =} \StringTok{"x"}\NormalTok{, }\AttributeTok{ylab =} \StringTok{"Densidad"}\NormalTok{)}
\end{Highlighting}
\end{Shaded}

\pandocbounded{\includegraphics[keepaspectratio]{Actividad-1_files/figure-latex/unnamed-chunk-5-1.pdf}}

\begin{Shaded}
\begin{Highlighting}[]
\CommentTok{\# Restaurar panel gráfico}
\FunctionTok{par}\NormalTok{(}\AttributeTok{mfrow =} \FunctionTok{c}\NormalTok{(}\DecValTok{1}\NormalTok{, }\DecValTok{1}\NormalTok{))}
\end{Highlighting}
\end{Shaded}

\begin{Shaded}
\begin{Highlighting}[]
\CommentTok{\# Configuración del panel para mostrar 3 gráficos}
\FunctionTok{par}\NormalTok{(}\AttributeTok{mfrow =} \FunctionTok{c}\NormalTok{(}\DecValTok{1}\NormalTok{, }\DecValTok{3}\NormalTok{))}

\CommentTok{\# Set de datos de ejemplo (n muestras)}
\NormalTok{n }\OtherTok{=} \DecValTok{10000}

\CommentTok{\# 1. Muy sesgada a la derecha: Exponencial (o Gamma alpha=1)}
\NormalTok{x21 }\OtherTok{=} \FunctionTok{rexp}\NormalTok{(n, }\AttributeTok{rate =} \DecValTok{1}\NormalTok{) }\CommentTok{\# equivalente a rgamma(n, shape = 1, rate = 1)}
\FunctionTok{hist}\NormalTok{(x21, }\AttributeTok{main =} \StringTok{"Muy sesgada a la derecha}\SpecialCharTok{\textbackslash{}n}\StringTok{(Exponencial lambda=1)"}\NormalTok{, }
     \AttributeTok{col =} \StringTok{"lightblue"}\NormalTok{, }\AttributeTok{xlab =} \StringTok{"x"}\NormalTok{, }\AttributeTok{freq =} \ConstantTok{FALSE}\NormalTok{)}

\CommentTok{\# 2. Ligeramente sesgada a la derecha: Weibull (shape=2)}
\NormalTok{x22 }\OtherTok{=} \FunctionTok{rweibull}\NormalTok{(n, }\AttributeTok{shape =} \DecValTok{2}\NormalTok{, }\AttributeTok{scale =} \DecValTok{1}\NormalTok{)}
\FunctionTok{hist}\NormalTok{(x22, }\AttributeTok{main =} \StringTok{"Ligeramente sesgada a la derecha}\SpecialCharTok{\textbackslash{}n}\StringTok{(Weibull shape=2)"}\NormalTok{, }
     \AttributeTok{col =} \StringTok{"lightgreen"}\NormalTok{, }\AttributeTok{xlab =} \StringTok{"x"}\NormalTok{, }\AttributeTok{freq =} \ConstantTok{FALSE}\NormalTok{)}

\CommentTok{\# 3. Sesgada a la izquierda: Weibull (shape=5)}
\NormalTok{x23 }\OtherTok{=} \FunctionTok{rweibull}\NormalTok{(n, }\AttributeTok{shape =} \DecValTok{10}\NormalTok{, }\AttributeTok{scale =} \DecValTok{1}\NormalTok{)}
\FunctionTok{hist}\NormalTok{(x23, }\AttributeTok{main =} \StringTok{"Sesgada a la izquierda}\SpecialCharTok{\textbackslash{}n}\StringTok{(Weibull shape=10)"}\NormalTok{, }
     \AttributeTok{col =} \StringTok{"coral"}\NormalTok{, }\AttributeTok{xlab =} \StringTok{"x"}\NormalTok{, }\AttributeTok{freq =} \ConstantTok{FALSE}\NormalTok{)}
\end{Highlighting}
\end{Shaded}

\pandocbounded{\includegraphics[keepaspectratio]{Actividad-1_files/figure-latex/unnamed-chunk-6-1.pdf}}

\begin{Shaded}
\begin{Highlighting}[]
\CommentTok{\# Restaurar panel gráfico}
\FunctionTok{par}\NormalTok{(}\AttributeTok{mfrow =} \FunctionTok{c}\NormalTok{(}\DecValTok{1}\NormalTok{, }\DecValTok{1}\NormalTok{))}
\end{Highlighting}
\end{Shaded}

\begin{Shaded}
\begin{Highlighting}[]
\CommentTok{\# 1. Muy sesgada a la derecha: Exponencial (o Gamma alpha=1)}
\NormalTok{sg21 }\OtherTok{=} \FunctionTok{skewness}\NormalTok{(x21)}
\NormalTok{k21 }\OtherTok{=} \FunctionTok{kurtosis}\NormalTok{(x21)}

\FunctionTok{cat}\NormalTok{(}\StringTok{"Sesgo:"}\NormalTok{, sg21, }\StringTok{"}\SpecialCharTok{\textbackslash{}n}\StringTok{"}\NormalTok{)}
\end{Highlighting}
\end{Shaded}

\begin{verbatim}
## Sesgo: 1.968455
\end{verbatim}

\begin{Shaded}
\begin{Highlighting}[]
\FunctionTok{cat}\NormalTok{(}\StringTok{"Curtosis:"}\NormalTok{, k21, }\StringTok{"}\SpecialCharTok{\textbackslash{}n}\StringTok{ }\SpecialCharTok{\textbackslash{}n}\StringTok{"}\NormalTok{)}
\end{Highlighting}
\end{Shaded}

\begin{verbatim}
## Curtosis: 8.725958 
## 
\end{verbatim}

\begin{Shaded}
\begin{Highlighting}[]
\CommentTok{\# 2. Ligeramente sesgada a la derecha: Weibull (shape=2)}
\NormalTok{sg22 }\OtherTok{=} \FunctionTok{skewness}\NormalTok{(x22)}
\NormalTok{k22 }\OtherTok{=} \FunctionTok{kurtosis}\NormalTok{(x22)}

\FunctionTok{cat}\NormalTok{(}\StringTok{"Sesgo:"}\NormalTok{, sg22, }\StringTok{"}\SpecialCharTok{\textbackslash{}n}\StringTok{"}\NormalTok{)}
\end{Highlighting}
\end{Shaded}

\begin{verbatim}
## Sesgo: 0.6243515
\end{verbatim}

\begin{Shaded}
\begin{Highlighting}[]
\FunctionTok{cat}\NormalTok{(}\StringTok{"Curtosis:"}\NormalTok{, k22, }\StringTok{"}\SpecialCharTok{\textbackslash{}n}\StringTok{ }\SpecialCharTok{\textbackslash{}n}\StringTok{"}\NormalTok{)}
\end{Highlighting}
\end{Shaded}

\begin{verbatim}
## Curtosis: 3.165923 
## 
\end{verbatim}

\begin{Shaded}
\begin{Highlighting}[]
\CommentTok{\# 3. Sesgada a la izquierda: Weibull (shape=5)}
\NormalTok{sg23 }\OtherTok{=} \FunctionTok{skewness}\NormalTok{(x23)}
\NormalTok{k23 }\OtherTok{=} \FunctionTok{kurtosis}\NormalTok{(x23)}

\FunctionTok{cat}\NormalTok{(}\StringTok{"Sesgo:"}\NormalTok{, sg23, }\StringTok{"}\SpecialCharTok{\textbackslash{}n}\StringTok{"}\NormalTok{)}
\end{Highlighting}
\end{Shaded}

\begin{verbatim}
## Sesgo: -0.6110021
\end{verbatim}

\begin{Shaded}
\begin{Highlighting}[]
\FunctionTok{cat}\NormalTok{(}\StringTok{"Curtosis:"}\NormalTok{, k23, }\StringTok{"}\SpecialCharTok{\textbackslash{}n}\StringTok{ }\SpecialCharTok{\textbackslash{}n}\StringTok{"}\NormalTok{)}
\end{Highlighting}
\end{Shaded}

\begin{verbatim}
## Curtosis: 3.541703 
## 
\end{verbatim}

\begin{Shaded}
\begin{Highlighting}[]
\CommentTok{\# Configuración del panel para mostrar 3 gráficos}
\FunctionTok{par}\NormalTok{(}\AttributeTok{mfrow =} \FunctionTok{c}\NormalTok{(}\DecValTok{1}\NormalTok{, }\DecValTok{3}\NormalTok{))}

\CommentTok{\#}
\FunctionTok{jarque.test}\NormalTok{(x21)}
\end{Highlighting}
\end{Shaded}

\begin{verbatim}
## 
##  Jarque-Bera Normality Test
## 
## data:  x21
## JB = 20119, p-value < 2.2e-16
## alternative hypothesis: greater
\end{verbatim}

\begin{Shaded}
\begin{Highlighting}[]
\FunctionTok{qqnorm}\NormalTok{(x21, }\AttributeTok{main =} \StringTok{"QQ{-}Plot Sesgada a la derecha"}\NormalTok{)}
\FunctionTok{qqline}\NormalTok{(x21, }\AttributeTok{col =} \StringTok{"red"}\NormalTok{)}

\FunctionTok{jarque.test}\NormalTok{(x22)}
\end{Highlighting}
\end{Shaded}

\begin{verbatim}
## 
##  Jarque-Bera Normality Test
## 
## data:  x22
## JB = 661.16, p-value < 2.2e-16
## alternative hypothesis: greater
\end{verbatim}

\begin{Shaded}
\begin{Highlighting}[]
\FunctionTok{qqnorm}\NormalTok{(x22, }\AttributeTok{main =} \StringTok{"QQ{-}Plot Ligeramente Sesgada a la derecha "}\NormalTok{)}
\FunctionTok{qqline}\NormalTok{(x22, }\AttributeTok{col =} \StringTok{"red"}\NormalTok{)}

\FunctionTok{jarque.test}\NormalTok{(x23)}
\end{Highlighting}
\end{Shaded}

\begin{verbatim}
## 
##  Jarque-Bera Normality Test
## 
## data:  x23
## JB = 744.47, p-value < 2.2e-16
## alternative hypothesis: greater
\end{verbatim}

\begin{Shaded}
\begin{Highlighting}[]
\FunctionTok{qqnorm}\NormalTok{(x23, }\AttributeTok{main =} \StringTok{"QQ{-}Plot Sesgada a la izquierda"}\NormalTok{)}
\FunctionTok{qqline}\NormalTok{(x23, }\AttributeTok{col =} \StringTok{"red"}\NormalTok{)}
\end{Highlighting}
\end{Shaded}

\pandocbounded{\includegraphics[keepaspectratio]{Actividad-1_files/figure-latex/unnamed-chunk-8-1.pdf}}

\begin{Shaded}
\begin{Highlighting}[]
\CommentTok{\# Restaurar panel gráfico}
\FunctionTok{par}\NormalTok{(}\AttributeTok{mfrow =} \FunctionTok{c}\NormalTok{(}\DecValTok{1}\NormalTok{, }\DecValTok{1}\NormalTok{))}
\end{Highlighting}
\end{Shaded}

\begin{itemize}
\tightlist
\item
  Compara los gráficos de la distribución teórica (punto A) con los
  gráficos de la distribución empírica (punto B), ¿en qué son iguales y
  en qué son diferentes?
\end{itemize}

Específicamente en las 3 distribuciones, tanto la teórica como la
empírica se ven casi idénticas: centradas en valores similares, con un
rango de variabilidad muy parecido y valor en las barras prácticamente
igual a lo esperado.

\begin{itemize}
\tightlist
\item
  Relaciona la distribución empírica (punto B) con los valores de sesgo
  y curtosis (punto C) y el QQPlot y la prueba de normalidad (punto D)
\end{itemize}

Nos iremos por cada distribución:

En el caso de la distribución muy sesgada a la derecha (exponencial),
observamos como su sesgo y curtosis son claramente mayores a 0, lo que
indica un sesgo prominente a la derecha así como una distribución
leptocúrtica. Mientras que la prueba de normalidad Jarque-Bera nos
indica que el p-value es muchísimo menor a 0.05, rechazando así la
hipótesis nula de normalidad; estos resultados son respaldador por la
QQ-Plot, donde la mayoría de los puntos no siguen la línea teórica y se
nota un claro sesgo a la derecha.

En la distribución ligeramente sesgada a la derecha (weibull shape=2)
podemos también notar valores de sesgo y curtosis superiores a 0, pero
no tan alejados como lo son los de la distribución exponencial; esto nos
indica lo mismo, sesgo a la derecha y distribución leptocúrtica. La
prueba de normalidad Jarque-Bera también rechaza la hipótesis nula y la
QQ-Plot, a pesar de tener una gran cantidad de valores siguiendo la
línea teórica, se observa diversos datos de ambas colas diferir con la
línea, sobretodo en la cola derecha.

Finalmente, tenemos la distribución sesgada a la izquierda (weibull
shape=10). En su sesgo, a diferencia de las distribuciones anteriores,
notamos que este es negativo, lo que nos indica efectivamente un sesgo a
la izquierda; la curtosis se mantiene por arriba de 0, por lo que
también es leptocúrtica. Con respecto a la prueba de normalidad, también
se rechaza la hipótesis nula y la QQ-Plot mostrada posee datos en ambas
colas que difieren con la línea teórica, mayormente en la cola
izquierda.

\begin{itemize}
\tightlist
\item
  Emite una conclusión de cómo se relaciona la curtosis y el sesgo con
  la distribución QQPlot y con la prueba de hipótesis de Jarque-Bera.
\end{itemize}

En conclusión, el sesgo y la curtosis representan de manera numérica la
forma y geometría de la distribución evaluada, mientras que estos
resultados se pueden ver de manera gráfica gracias a la QQ-Plot. En el
caso de la prueba de hipótesis de Jarque-Bera, esta misma combina tanto
el sesgo como la curtosis para así armar un método que nos ayude a
evaluar la normalidad de una distribución. Funciona de manera similar a
una prueba de hipótesis, donde en este caso la hipótesis nula es la
normalidad y la hipótesis alternativa es lo contrario.

\subsection{Probando normalidad en una base de
datos}\label{probando-normalidad-en-una-base-de-datos}

\begin{Shaded}
\begin{Highlighting}[]
\FunctionTok{library}\NormalTok{(nortest)}
\FunctionTok{library}\NormalTok{(e1071)}
\end{Highlighting}
\end{Shaded}

\begin{verbatim}
## 
## Adjuntando el paquete: 'e1071'
\end{verbatim}

\begin{verbatim}
## The following objects are masked from 'package:moments':
## 
##     kurtosis, moment, skewness
\end{verbatim}

\begin{Shaded}
\begin{Highlighting}[]
\FunctionTok{data}\NormalTok{(cars)}
\NormalTok{speed }\OtherTok{=}\NormalTok{ cars}\SpecialCharTok{$}\NormalTok{speed}
\NormalTok{dist }\OtherTok{=}\NormalTok{ cars}\SpecialCharTok{$}\NormalTok{dist}
\end{Highlighting}
\end{Shaded}

\begin{Shaded}
\begin{Highlighting}[]
\CommentTok{\# Prueba para \textquotesingle{}speed\textquotesingle{}}
\FunctionTok{shapiro.test}\NormalTok{(speed)}
\end{Highlighting}
\end{Shaded}

\begin{verbatim}
## 
##  Shapiro-Wilk normality test
## 
## data:  speed
## W = 0.97765, p-value = 0.4576
\end{verbatim}

\begin{Shaded}
\begin{Highlighting}[]
\FunctionTok{ad.test}\NormalTok{(speed)}
\end{Highlighting}
\end{Shaded}

\begin{verbatim}
## 
##  Anderson-Darling normality test
## 
## data:  speed
## A = 0.26143, p-value = 0.6927
\end{verbatim}

\begin{Shaded}
\begin{Highlighting}[]
\CommentTok{\# Prueba para \textquotesingle{}dist\textquotesingle{}}
\FunctionTok{shapiro.test}\NormalTok{(dist)}
\end{Highlighting}
\end{Shaded}

\begin{verbatim}
## 
##  Shapiro-Wilk normality test
## 
## data:  dist
## W = 0.95144, p-value = 0.0391
\end{verbatim}

\begin{Shaded}
\begin{Highlighting}[]
\FunctionTok{ad.test}\NormalTok{(dist)}
\end{Highlighting}
\end{Shaded}

\begin{verbatim}
## 
##  Anderson-Darling normality test
## 
## data:  dist
## A = 0.74067, p-value = 0.05021
\end{verbatim}

Se eligieron las pruebas de normalidad de Shapiro Wilk y Anderson
Darling, donde ambas variables fueron sometidas a estos dos métodos.
Ante estas pruebas de normalidad, tenemos lo siguiente:

\begin{itemize}
\tightlist
\item
  h0: La variable se distribuye de manera normal.
\item
  h1: La variable no se distribuye de manera normal.
\end{itemize}

En la variable speed, se muestran como ambos valores superan el p-value
= 0.05, lo que nos indica que la h0 se mantiene, así que speed sigue una
distribución normal. Por otro lado, en la variable dist ambas pruebas
salen con un p-value menor a 0.05, lo que refuta la hipótesis nula y nos
indica que esta variable no sigue una distribución normal.

\begin{Shaded}
\begin{Highlighting}[]
\FunctionTok{par}\NormalTok{(}\AttributeTok{mfrow =} \FunctionTok{c}\NormalTok{(}\DecValTok{1}\NormalTok{, }\DecValTok{2}\NormalTok{))}

\FunctionTok{qqnorm}\NormalTok{(speed, }\AttributeTok{main =} \StringTok{"QQ{-}Plot Speed"}\NormalTok{)}
\FunctionTok{qqline}\NormalTok{(speed, }\AttributeTok{col =} \StringTok{"red"}\NormalTok{, }\AttributeTok{lwd =} \DecValTok{2}\NormalTok{)}

\FunctionTok{qqnorm}\NormalTok{(dist, }\AttributeTok{main =} \StringTok{"QQ{-}Plot Dist"}\NormalTok{)}
\FunctionTok{qqline}\NormalTok{(dist, }\AttributeTok{col =} \StringTok{"red"}\NormalTok{, }\AttributeTok{lwd =} \DecValTok{2}\NormalTok{)}
\end{Highlighting}
\end{Shaded}

\pandocbounded{\includegraphics[keepaspectratio]{Actividad-1_files/figure-latex/unnamed-chunk-11-1.pdf}}

En la gráfica de speed se observa que la mayor parte de los puntos
siguen la recta teórica de normalidad. Mientras que en la variable dist,
se observa desfase en ambas colas, especialmente en la de la derecha.

\begin{Shaded}
\begin{Highlighting}[]
\CommentTok{\# Sesgo y Curtosis para \textquotesingle{}speed\textquotesingle{}}
\FunctionTok{skewness}\NormalTok{(speed)}
\end{Highlighting}
\end{Shaded}

\begin{verbatim}
## [1] -0.1105533
\end{verbatim}

\begin{Shaded}
\begin{Highlighting}[]
\FunctionTok{kurtosis}\NormalTok{(speed)}
\end{Highlighting}
\end{Shaded}

\begin{verbatim}
## [1] -0.6730924
\end{verbatim}

\begin{Shaded}
\begin{Highlighting}[]
\CommentTok{\# Sesgo y Curtosis para \textquotesingle{}dist\textquotesingle{}}
\FunctionTok{skewness}\NormalTok{(dist)}
\end{Highlighting}
\end{Shaded}

\begin{verbatim}
## [1] 0.7591268
\end{verbatim}

\begin{Shaded}
\begin{Highlighting}[]
\FunctionTok{kurtosis}\NormalTok{(dist)}
\end{Highlighting}
\end{Shaded}

\begin{verbatim}
## [1] 0.1193971
\end{verbatim}

En el caso de la variable speed, se tiene una simetría cercana a 0, lo
que indica efectivamente simetría, mientras que al ser su curtosis menor
a 0, el resultado nos indica que la distribución es platicúrtica. En la
variable dist, se observa como el valor de sesgo está más alejado del 0,
y al ser positivo nos indica un ligero sesgo a la derecha; mientras que
en su curtosis parece ser leptocúrtica al ser mayor a 0.

\begin{Shaded}
\begin{Highlighting}[]
\CommentTok{\# Para \textquotesingle{}speed\textquotesingle{}}
\NormalTok{mean\_speed }\OtherTok{\textless{}{-}} \FunctionTok{mean}\NormalTok{(speed)}
\NormalTok{median\_speed }\OtherTok{\textless{}{-}} \FunctionTok{median}\NormalTok{(speed)}
\NormalTok{midrange\_speed }\OtherTok{\textless{}{-}}\NormalTok{ (}\FunctionTok{max}\NormalTok{(speed) }\SpecialCharTok{+} \FunctionTok{min}\NormalTok{(speed)) }\SpecialCharTok{/} \DecValTok{2}

\FunctionTok{cat}\NormalTok{(}\StringTok{"SPEED {-} Media:"}\NormalTok{, mean\_speed, }\StringTok{"| Mediana:"}\NormalTok{, median\_speed, }\StringTok{"| Rango Medio:"}\NormalTok{, midrange\_speed, }\StringTok{"}\SpecialCharTok{\textbackslash{}n}\StringTok{"}\NormalTok{)}
\end{Highlighting}
\end{Shaded}

\begin{verbatim}
## SPEED - Media: 15.4 | Mediana: 15 | Rango Medio: 14.5
\end{verbatim}

\begin{Shaded}
\begin{Highlighting}[]
\CommentTok{\# Para \textquotesingle{}dist\textquotesingle{}}
\NormalTok{mean\_dist }\OtherTok{\textless{}{-}} \FunctionTok{mean}\NormalTok{(dist)}
\NormalTok{median\_dist }\OtherTok{\textless{}{-}} \FunctionTok{median}\NormalTok{(dist)}
\NormalTok{midrange\_dist }\OtherTok{\textless{}{-}}\NormalTok{ (}\FunctionTok{max}\NormalTok{(dist) }\SpecialCharTok{+} \FunctionTok{min}\NormalTok{(dist)) }\SpecialCharTok{/} \DecValTok{2}

\FunctionTok{cat}\NormalTok{(}\StringTok{"DIST {-} Media:"}\NormalTok{, mean\_dist, }\StringTok{"| Mediana:"}\NormalTok{, median\_dist, }\StringTok{"| Rango Medio:"}\NormalTok{, midrange\_dist, }\StringTok{"}\SpecialCharTok{\textbackslash{}n}\StringTok{"}\NormalTok{)}
\end{Highlighting}
\end{Shaded}

\begin{verbatim}
## DIST - Media: 42.98 | Mediana: 36 | Rango Medio: 61
\end{verbatim}

La distancia entre las tres medidas en la variable speed es mínima, lo
que nos indica simetría. En cambio, se observa una diferencia
considerable entre las 3 medidas de dist, sobretodo en el rango medio,
lo que nos da una idea de cómo es su verdadera distribución.

\begin{Shaded}
\begin{Highlighting}[]
\FunctionTok{par}\NormalTok{(}\AttributeTok{mfrow =} \FunctionTok{c}\NormalTok{(}\DecValTok{1}\NormalTok{, }\DecValTok{2}\NormalTok{))}

\FunctionTok{boxplot}\NormalTok{(speed, }\AttributeTok{main =} \StringTok{"Boxplot Velocidad"}\NormalTok{, }\AttributeTok{ylab =} \StringTok{"mph"}\NormalTok{, }\AttributeTok{col =} \StringTok{"lightblue"}\NormalTok{)}
\FunctionTok{boxplot}\NormalTok{(dist, }\AttributeTok{main =} \StringTok{"Boxplot Distancia"}\NormalTok{, }\AttributeTok{ylab =} \StringTok{"ft"}\NormalTok{, }\AttributeTok{col =} \StringTok{"lightgreen"}\NormalTok{)}
\end{Highlighting}
\end{Shaded}

\pandocbounded{\includegraphics[keepaspectratio]{Actividad-1_files/figure-latex/unnamed-chunk-14-1.pdf}}

La boxplot de la variable sesgo parece estar distribuida de manera
normal, al no tener datos atípicos y mostrar al rango medio y a la
mediana muy cercanas. En cambio, en el boxplot de dist se observa
claramente el sesgo que posee a la derecha, donde podemos observar que
gran parte de los datos se encuentran a la izquierda pero hay ciertos
datos atípicos que van a la derecha y afectan a la media (de hecho, en
el boxplot se puede observar un dato atípico muy prominente).

\begin{Shaded}
\begin{Highlighting}[]
\FunctionTok{par}\NormalTok{(}\AttributeTok{mfrow =} \FunctionTok{c}\NormalTok{(}\DecValTok{1}\NormalTok{, }\DecValTok{2}\NormalTok{))}

\CommentTok{\# Histograma para \textquotesingle{}speed\textquotesingle{}}
\FunctionTok{hist}\NormalTok{(speed, }\AttributeTok{freq =} \ConstantTok{FALSE}\NormalTok{, }\AttributeTok{main =} \StringTok{"Histograma Velocidad"}\NormalTok{, }\AttributeTok{xlab =} \StringTok{"Speed"}\NormalTok{, }\AttributeTok{col =} \StringTok{"lightblue"}\NormalTok{)}
\FunctionTok{lines}\NormalTok{(}\FunctionTok{density}\NormalTok{(speed), }\AttributeTok{col =} \StringTok{"red"}\NormalTok{, }\AttributeTok{lwd =} \DecValTok{2}\NormalTok{)}
\FunctionTok{curve}\NormalTok{(}\FunctionTok{dnorm}\NormalTok{(x, }\AttributeTok{mean =} \FunctionTok{mean}\NormalTok{(speed), }\AttributeTok{sd =} \FunctionTok{sd}\NormalTok{(speed)), }
      \AttributeTok{from =} \FunctionTok{min}\NormalTok{(speed), }\AttributeTok{to =} \FunctionTok{max}\NormalTok{(speed), }\AttributeTok{add =} \ConstantTok{TRUE}\NormalTok{, }\AttributeTok{col =} \StringTok{"blue"}\NormalTok{, }\AttributeTok{lwd =} \DecValTok{2}\NormalTok{)}

\CommentTok{\# Histograma para \textquotesingle{}dist\textquotesingle{}}
\FunctionTok{hist}\NormalTok{(dist, }\AttributeTok{freq =} \ConstantTok{FALSE}\NormalTok{, }\AttributeTok{main =} \StringTok{"Histograma Distancia"}\NormalTok{, }\AttributeTok{xlab =} \StringTok{"Dist"}\NormalTok{, }\AttributeTok{col =} \StringTok{"lightgreen"}\NormalTok{)}
\FunctionTok{lines}\NormalTok{(}\FunctionTok{density}\NormalTok{(dist), }\AttributeTok{col =} \StringTok{"red"}\NormalTok{, }\AttributeTok{lwd =} \DecValTok{2}\NormalTok{)}
\FunctionTok{curve}\NormalTok{(}\FunctionTok{dnorm}\NormalTok{(x, }\AttributeTok{mean =} \FunctionTok{mean}\NormalTok{(dist), }\AttributeTok{sd =} \FunctionTok{sd}\NormalTok{(dist)), }
      \AttributeTok{from =} \FunctionTok{min}\NormalTok{(dist), }\AttributeTok{to =} \FunctionTok{max}\NormalTok{(dist), }\AttributeTok{add =} \ConstantTok{TRUE}\NormalTok{, }\AttributeTok{col =} \StringTok{"blue"}\NormalTok{, }\AttributeTok{lwd =} \DecValTok{2}\NormalTok{)}
\end{Highlighting}
\end{Shaded}

\pandocbounded{\includegraphics[keepaspectratio]{Actividad-1_files/figure-latex/unnamed-chunk-15-1.pdf}}

\begin{Shaded}
\begin{Highlighting}[]
\FunctionTok{par}\NormalTok{(}\AttributeTok{mfrow =} \FunctionTok{c}\NormalTok{(}\DecValTok{1}\NormalTok{, }\DecValTok{1}\NormalTok{))}
\end{Highlighting}
\end{Shaded}

En la variable speed, se muestra como la curva de densidad empírica
encaja en su mayoría con la curva normal teórica, a excepción de que la
curva primero mencionada se muestra más platicúrtica. En el caso de la
variable dist, la curva de densidad empírica muestra un claro sesgo a la
derecha, así como también es más puntiaguda que la curva de la normal
teórica al ser leptocúrtica.

\begin{itemize}
\tightlist
\item
  Emite una conclusión final sobre la normalidad de los datos en
  términos de todas las pruebas que realizaste.
\end{itemize}

Finalmente, por todas las pruebas mostradas, concluimos que la
distribución de la variable speed sí cumple con gran parte de las
pruebas de normalidad, por lo que podemos afirmar la hipótesis nula.
Mientras que la distribución de dist no cumplió de manera satisfactoria
con ninguna de las pruebas, por lo que se rechaza la hipótesis nula de
normalidad.

\subsection{Declaratoria de uso de IA durante el desarrollo del
trabajo}\label{declaratoria-de-uso-de-ia-durante-el-desarrollo-del-trabajo}

Se utilizó IA:

\begin{enumerate}
\def\labelenumi{(\arabic{enumi})}
\item
  Herramienta: Gemini
\item
  Uso realizado: Ayuda en el código al momento de importar la librería
  cars, generación de las distribuciones de la segunda mitad de la
  primera sección, y sintaxis para la realización de elementos visuales
  como histogramas, boxplots, curvas de densidad entre otras.
\item
  Secciones donde se utilizó: Segunda mitad de la primera sección,
  específicamente en los incisos A y B, y a lo largo de la última
  sección en la importación de la librería, las gráficas QQ-Plot, las
  boxplots y los histogramas.
\item
  Validación realizada por los estudiantes: Se verificó, corrió y depuró
  todo el script de R, así como también analizaron todos los resultados
  arrojados para validar que todo estuviera correcto y en orden.
\item
  Confirmo de que se revisó, verificó y que se asume la responsabilidad
  del contenido final entregado.
\end{enumerate}

\end{document}
